\section{失效模式谱系：从对抗扰动到策略奶酪}

讲者用游戏术语把失败分成三个层次：adversarial policies、exploiters、cheese。这一分类对 LLM 安全评估极具操作性，因为它对应了不同强度、不同成本的攻击者模型。

\begin{knowledgebox}{三层失效模型}
\begin{itemize}
    \item \textbf{Adversarial-like}：微小扰动触发大幅错误，类似对抗样本。
    \item \textbf{Exploiter}：持续围绕某一弱点迭代，形成稳定压制。
    \item \textbf{Cheese}：策略可泛化到多个对手，但不一定是整体最优。
\end{itemize}
\end{knowledgebox}

\begin{warningbox}{为什么 ``How many r in strawberry'' 被反复提及}
讲者用这个例子强调：系统可能在高难任务表现亮眼，却在基础一致性上失败。这并非``偶发小错''，而是能力分布不均衡，提示我们要做 failure taxonomy，而非只看均值指标。
\end{warningbox}

\begin{importantbox}{从失败分层到测试分层}
对应三层失效模型，测试也要分层：
\begin{itemize}
    \item 扰动稳健性测试（格式噪声、提示变体、边界输入）
    \item 目标导向 exploit 测试（反复 probing、记忆污染、工具串联攻击）
    \item 分布迁移测试（新任务族、长时交互、跨工具组合）
\end{itemize}
\end{importantbox}

这部分最重要的一句话是：\textit{``You have these models, they do amazing things, but we see them fail often.''}\footnote{来源：\href{https://www.youtube.com/watch?v=ntjOxjZMaac}{Lecture 03 字幕}，区间 00:32:57--00:33:03。} 这不是悲观论断，而是训练目标需要升级的证据。

\subsection{本章小结}
失败不是单点事件，而是一个分层谱系。只有将其结构化，才能把鲁棒性从``运气''变成``可优化对象''，这也是后续 league 机制存在的根本理由。

